\documentclass[UTF8,a4paper,12pt]{ctexart}

\usepackage{amsmath,amssymb,bm}
\usepackage{geometry}
\usepackage{graphicx}
\usepackage{booktabs}
\usepackage{physics}
\usepackage{hyperref}
\usepackage{subcaption}
\usepackage{cite}

\geometry{left=2.5cm,right=2.5cm,top=2.5cm,bottom=2.5cm}

\title{二维薛定谔方程的谱方法数值推导}
\author{王博宇}
\date{}

\begin{document}

\maketitle

\section{问题背景}

薛定谔方程是量子力学中的基本控制方程，用于描述微观粒子波函数随时间的演化规律
\cite{griffiths,sakurai}。
在经典力学中，粒子的状态通常由位置与速度决定；
而在量子力学中，粒子的状态由复值波函数
$\psi(x,y,t)$
决定，其模平方
$|\psi(x,y,t)|^2$
表示粒子在空间中的概率密度分布。

时变薛定谔方程（Time-Dependent Schrödinger Equation, TDSE）具有重要的理论与工程意义。例如：
电子在势场中的传播、原子与分子的量子态演化、量子隧穿与干涉现象、
半导体器件中的电子输运、光波与量子波包传播模拟，以及 Bose--Einstein 凝聚态系统都属于其典型应用场景。

然而，薛定谔方程通常难以获得解析解。
除了少量理想模型（如无限深势阱、简谐振子、氢原子等）外，
大多数实际问题都需要依赖数值方法进行求解。

本文考虑二维时变薛定谔方程：

\[
i\frac{\partial \psi}{\partial t}
=
-\Delta \psi + V(x,y)\psi,
\]

其中 $ \psi(x,y,t) $ 为二维复值波函数；$
    \Delta
    =
    \frac{\partial^2}{\partial x^2}
    +
    \frac{\partial^2}{\partial y^2} $
为二维拉普拉斯算子；
$ V(x,y) $
表示外部势场。

其中拉普拉斯项对应粒子的动能项，而势场项描述外部环境对粒子的影响。

本文希望解决的问题包括：
如何数值模拟二维波函数的传播过程，如何稳定计算波函数的空间二阶导数，
如何高效求解大规模二维网格上的薛定谔方程，如何直观展示复值波函数的相位与传播行为，
以及比较不同数值方法（有限差分与 FFT 谱方法）的性能差异。

在传统有限差分方法中，拉普拉斯算子通常采用局域差分近似：

\[
\Delta \psi_{i,j}
\approx
\psi_{i+1,j}
+
\psi_{i-1,j}
+
\psi_{i,j+1}
+
\psi_{i,j-1}
-
4\psi_{i,j}.
\]

这种方法实现简单，但会带来较明显的数值色散误差，
并且在高频传播问题中容易产生不稳定现象。

为了获得更高精度，本文采用傅里叶谱方法（Fourier Spectral Method）
计算拉普拉斯算符 \cite{trefethen,boyd}。
其核心思想是利用傅里叶变换将微分运算转化为频域中的代数乘法。

设：

\[
\hat\psi(k_x,k_y)
=
\mathcal{F}[\psi(x,y)],
\]

则有：

\[
\mathcal{F}[\Delta\psi]
=
-(k_x^2+k_y^2)\hat\psi.
\]

因此空间二阶导数可以通过 FFT 快速计算，而无需显式构造差分模板。

快速傅里叶变换（Fast Fourier Transform, FFT）
可以高效实现频域计算 \cite{fft}。
相比传统有限差分方法，FFT 谱方法具有：
更高的空间精度、更低的数值色散、更平滑的波传播效果，
并且更适合处理周期性与波动问题。

因此，FFT 谱方法已经成为求解时变薛定谔方程、
Gross--Pitaevskii 方程以及各种波动方程的重要数值工具。

本文将基于 Python、NumPy 与 Pygame，
实现二维量子波包传播的可视化模拟，
并比较有限差分方法与 FFT 谱方法在数值表现与运行效率上的差异。

\section{方法提出}

本文研究二维时变薛定谔方程的数值求解与可视化问题。为了模拟量子波包在二维空间中的传播行为，本文采用傅里叶谱方法（Fourier Spectral Method）对空间导数进行离散，并结合显式时间推进格式实现波函数演化。

\subsection{二维薛定谔方程}

本文考虑二维时变薛定谔方程：

\[
i\frac{\partial \psi}{\partial t}
=
-\Delta \psi + V(x,y)\psi,
\]

其中：
$\psi(x,y,t)$ 为二维复值波函数；
$\Delta=\frac{\partial^2}{\partial x^2}+\frac{\partial^2}{\partial y^2}$ 为二维拉普拉斯算子；
$V(x,y)$ 表示外部势场。

其中拉普拉斯项对应粒子的动能项，而势场项描述外部环境对量子粒子的作用。

\subsection{初始波包构造}

为了模拟局域化量子粒子的传播，本文采用二维高斯波包作为初始条件：

\[
\psi(x,y,0)
=
\exp\left(
-\frac{(x-x_0)^2+(y-y_0)^2}{2\sigma^2}
\right)
\exp\left(
i2\pi p(y-y_0)
\right).
\]

其中：
第一项为高斯包络，用于控制粒子的空间局域性；
$(x_0,y_0)$ 为初始中心位置；$\sigma$ 控制波包宽度；
第二项为平面波相位，用于赋予波包初始动量；$p$ 表示初始波数大小。

因此，该初始条件可以理解为“具有初始传播方向的局域化量子粒子”。

\subsection{波函数归一化}

量子力学中，波函数模平方表示概率密度： $\rho(x,y,t)=|\psi(x,y,t)|^2.$，其中概率总和必须满足： $\iint_\Omega |\psi(x,y,t)|^2 \dd x \dd y = 1.$，
离散化后，设网格尺寸为：
$ \Delta x=\frac1{N_x},
\qquad
\Delta y=\frac1{N_y},
$ 则归一化因子为：
$
C
=
\sum_{j,k}
|\psi_{j,k}|^2
\Delta x\Delta y.
$
因此每一步演化后需要执行：
$
\psi_{j,k}
\leftarrow
\frac{\psi_{j,k}}{\sqrt{C}}.
$

由于本文采用显式 Euler 时间推进，其并不严格满足幺正性，因此重新归一化能够有效抑制数值误差累积。

\subsection{傅里叶谱方法}

传统有限差分方法通常采用局域模板近似拉普拉斯算子：

\[
\Delta \psi_{i,j}
\approx
\psi_{i+1,j}
+
\psi_{i-1,j}
+
\psi_{i,j+1}
+
\psi_{i,j-1}
-
4\psi_{i,j}.
\]

虽然实现简单，但会带来明显数值色散误差，并且在高频波传播问题中容易产生数值不稳定。

因此本文采用傅里叶谱方法计算空间导数
\cite{trefethen,boyd}。

设：

\[
\hat\psi(k_x,k_y)
=
\mathcal{F}[\psi(x,y)],
\]

根据傅里叶变换微分性质，有：

\[
\mathcal{F}
\left[
\frac{\partial^2\psi}{\partial x^2}
\right]
=
-k_x^2\hat\psi,
\]

\[
\mathcal{F}
\left[
\frac{\partial^2\psi}{\partial y^2}
\right]
=
-k_y^2\hat\psi.
\]

因此：

\[
\mathcal{F}[\Delta\psi]
=
-(k_x^2+k_y^2)\hat\psi.
\]

记：

$k^2=k_x^2+k_y^2,$

则：

\[
\widehat{\Delta\psi}
=
-k^2\hat\psi.
\]

因此可以通过 FFT 快速计算拉普拉斯项：

\[
\Delta\psi
=
\mathcal{F}^{-1}
\left[
-k^2\mathcal{F}[\psi]
\right].
\]

快速傅里叶变换（FFT）能够高效实现频域计算
\cite{fft}。

相比有限差分方法，傅里叶谱方法具有：
更高空间精度、更低数值色散、更平滑波前传播，
并且更适合波动问题。

\subsection{哈密顿算符离散化}

薛定谔方程中的哈密顿算符定义为：

\[
\hat H\psi
=
-\Delta\psi+V\psi.
\]

利用傅里叶谱方法，可以得到：

\[
-\Delta\psi
=
\mathcal{F}^{-1}
\left[
k^2\hat\psi
\right].
\]

势场项直接在实空间计算：

$V\psi=V(x,y)\psi(x,y).$

因此哈密顿算符作用结果为：

\[
\hat H\psi
=
\mathcal{F}^{-1}
\left[
k^2\hat\psi
\right]
+
V\psi.
\]

\subsection{时间推进格式}

由薛定谔方程：

\[
\frac{\partial\psi}{\partial t}
=
-i\hat H\psi,
\]

本文采用显式 Euler 时间推进：

\[
\psi^{n+1}
=
\psi^n
+
\Delta t
\left(
-i\hat H\psi^n
\right).
\]

代入哈密顿算符后得到：

\[
\psi^{n+1}
=
\psi^n
-i\Delta t
\left[
\mathcal{F}^{-1}
\left(
k^2\mathcal{F}[\psi^n]
\right)
+
V\psi^n
\right].
\]

该公式即为程序中每一步波函数更新的核心形式。

\subsection{边界条件}

为了模拟有限空间区域中的粒子传播，本文在边界处施加零边界条件：

$\psi(0,y,t)=0,$
$\psi(1,y,t)=0,$
$\psi(x,0,t)=0,$
$\psi(x,1,t)=0.$

该条件等价于无限高势垒边界，即粒子无法离开计算区域。

\subsection{波包统计量}

为了分析波函数传播行为，本文进一步计算粒子位置期望值：

\[
\mathbb E[x]
=
\iint x|\psi|^2 \dd x\dd y,
\]

\[
\mathbb E[y]
=
\iint y|\psi|^2 \dd x\dd y.
\]

定义波包中心：

\[
\bm r_c
=
(\mathbb E[x],\mathbb E[y]).
\]

同时定义波包扩散尺度：

\[
D_r
=
\iint
\left[
(x-\mathbb E[x])^2
+
(y-\mathbb E[y])^2
\right]
|\psi|^2
\dd x\dd y.
\]

其平方根：

$\sigma_r=\sqrt{D_r}$

用于描述波包扩散程度。

\subsection{势场模型}

本文采用中心势场：

\[
V(x,y)
=
-\frac{\alpha}
{
(x-1/2)^2+(y-1/2)^2+\varepsilon
}.
\]

其中：
$\alpha>0$ 时为吸引势，$\alpha<0$ 时为排斥势，$\varepsilon$ 用于避免分母奇异。

当：

$\alpha=0$

时，系统退化为自由粒子传播模型。

\subsection{方法总结}

综合上述过程，本文提出了一种基于 FFT 谱方法的二维薛定谔方程数值求解框架。

整体流程如下：
首先构造高斯波包初始条件并对波函数进行归一化，然后利用 FFT 计算空间拉普拉斯项并构造哈密顿算符，
接着使用显式 Euler 方法推进时间，最后施加边界条件并重新归一化，同时可视化波函数传播过程。

该方法能够较高效地模拟二维量子波包在势场中的传播、扩散与干涉行为。

\section{数值实验}

为了验证本文提出的二维薛定谔方程数值求解方法，本文分别进行了：
基于 FFT 谱方法的二维量子波包传播实验，以及有限差分方法与 FFT 谱方法的对比实验。

实验主要观察波函数在二维空间中的传播、扩散、干涉以及边界反射行为，并分析不同数值方法的计算效率与数值稳定性。

\subsection{实验设置}

实验采用二维高斯波包作为初始条件：

\[
\psi(x,y,0)
=
\exp\left(
-\frac{(x-x_0)^2+(y-y_0)^2}{2\sigma^2}
\right)
\exp\left(
i2\pi p(y-y_0)
\right).
\]

其中：

\[
x_0=0.5,
\qquad
y_0=0.2,
\]

波包宽度为：

$\sigma=0.1,$

初始波数为：

$p=10.$

实验采用二维离散网格：

$512\times512,$

时间步长设置为：

$dt=10^{-3}.$

程序在每一个时间步均保存波函数图像，用于后续分析与视频生成。

\subsection{FFT谱方法传播结果}

图~\ref{fig:wave_packet_series} 展示了基于 FFT 谱方法得到的二维波函数传播结果。
\begin{figure}[htbp]
    \centering

    % ================= FFT =================
    \begin{subfigure}{0.16\linewidth}
        \centering
        \includegraphics[width=\linewidth]{../field_state/electron_img/frame_0000.jpg}
        \caption{$t=0$}
    \end{subfigure}
    \hfill
    \begin{subfigure}{0.16\linewidth}
        \centering
        \includegraphics[width=\linewidth]{../field_state/electron_img/frame_0300.jpg}
        \caption{$t=300$}
    \end{subfigure}
    \hfill
    \begin{subfigure}{0.16\linewidth}
        \centering
        \includegraphics[width=\linewidth]{../field_state/electron_img/frame_0600.jpg}
        \caption{$t=600$}
    \end{subfigure}
    \hfill
    \begin{subfigure}{0.16\linewidth}
        \centering
        \includegraphics[width=\linewidth]{../field_state/electron_img/frame_0900.jpg}
        \caption{$t=900$}
    \end{subfigure}
    \hfill
    \begin{subfigure}{0.16\linewidth}
        \centering
        \includegraphics[width=\linewidth]{../field_state/electron_img/frame_1200.jpg}
        \caption{$t=1200$}
    \end{subfigure}
    \hfill
    \begin{subfigure}{0.16\linewidth}
        \centering
        \includegraphics[width=\linewidth]{../field_state/electron_img/frame_1500.jpg}
        \caption{$t=1500$}
    \end{subfigure}

    \caption{FFT谱方法下，拥有初始动量的二维自由电子波包传播过程}
    \label{fig:wave_packet_series}

    \vspace{0.3cm}

    % ================= Euler =================
    \begin{subfigure}{0.16\linewidth}
        \centering
        \includegraphics[width=\linewidth]{../field_state/electron_img_euler/frame_0000.jpg}
        \caption{$t=0$}
    \end{subfigure}
    \hfill
    \begin{subfigure}{0.16\linewidth}
        \centering
        \includegraphics[width=\linewidth]{../field_state/electron_img_euler/frame_0300.jpg}
        \caption{$t=300$}
    \end{subfigure}
    \hfill
    \begin{subfigure}{0.16\linewidth}
        \centering
        \includegraphics[width=\linewidth]{../field_state/electron_img_euler/frame_0600.jpg}
        \caption{$t=600$}
    \end{subfigure}
    \hfill
    \begin{subfigure}{0.16\linewidth}
        \centering
        \includegraphics[width=\linewidth]{../field_state/electron_img_euler/frame_0900.jpg}
        \caption{$t=900$}
    \end{subfigure}
    \hfill
    \begin{subfigure}{0.16\linewidth}
        \centering
        \includegraphics[width=\linewidth]{../field_state/electron_img_euler/frame_1200.jpg}
        \caption{$t=1200$}
    \end{subfigure}
    \hfill
    \begin{subfigure}{0.16\linewidth}
        \centering
        \includegraphics[width=\linewidth]{../field_state/electron_img_euler/frame_1500.jpg}
        \caption{$t=1500$}
    \end{subfigure}

    \caption{有限差分方法下二维量子波包传播过程。可以看到在相同时间步长下，有限差分方法更容易出现局部振荡与数值误差，而 FFT 谱方法能够在更细网格下保持较稳定传播。}
    \label{fig:euler_series}

\end{figure}

程序并未直接显示概率密度：

$|\psi(x,y,t)|^2,$

而是对波函数实部与虚部分别进行颜色编码。

设：

\[
\psi(x,y,t)
=
\Re(\psi)+i\Im(\psi),
\]

则：
红色通道表示波函数实部，蓝色通道表示波函数虚部，黑色区域表示波函数幅值较小。

因此：
红色越亮表示实部越大，蓝色越亮表示虚部越大，紫色区域表示实部与虚部同时较强，黑色区域表示波函数接近零。

相比单独显示概率密度，这种复平面颜色编码能够更直观展示波函数相位旋转与干涉结构。

从图中可以观察到：
初始高斯波包沿 $y$ 方向传播，波函数相位不断旋转并导致颜色持续变化，
波包在传播过程中逐渐扩散，边界附近出现反射与干涉现象，
且波函数整体保持连续和平滑。

这些现象符合量子波传播的基本规律。

\subsection{有限差分方法对比实验}

为了比较不同数值方法的表现，本文进一步实现了基于有限差分拉普拉斯算子的显式时间推进方法。

有限差分拉普拉斯采用五点格式：

\[
\Delta \psi_{i,j}
=
\psi_{i+1,j}
+
\psi_{i-1,j}
+
\psi_{i,j+1}
+
\psi_{i,j-1}
-
4\psi_{i,j}.
\]

程序实现如下：

\begin{verbatim}
lap = np.roll(psi,1,axis=0)
    + np.roll(psi,-1,axis=0)
    + np.roll(psi,1,axis=1)
    + np.roll(psi,-1,axis=1)
    - 4*psi
\end{verbatim}

同时，为提高时间精度，有限差分方法采用二阶 Taylor 展开：

\[
\psi(t+\Delta t)
=
\psi(t)
-
iH\psi \Delta t
-
\frac12 H^2\psi \Delta t^2.
\]

其中：

$H\psi=-\Delta\psi+V\psi.$

图~\ref{fig:euler_series} 展示了有限差分方法下的传播结果。


从图中可以观察到：
波包同样会发生扩散，相位不断变化，边界附近会出现明显数值反射，
且局部高频振荡更加明显。

相比之下，FFT 谱方法得到的波前更加连续和平滑。

\subsection{计算效率分析}

有限差分方法的单步复杂度约为：
$O(N),$，
其中 $N$ 为网格点数量。
然而由于显式差分方法稳定性较差，需要满足：

$dt < (\Delta x)^2,$

因此往往需要极小时间步长。
另一方面，FFT 谱方法虽然单步需要进行二维 FFT：

$O(N\log N),$

但其空间精度更高，数值色散更小，因此能够在较低网格分辨率下获得更平滑结果。

实验中可以明显观察到：
FFT 方法整体传播更加稳定，波前更加连续，并且在较低分辨率下即可获得较好结果；
有限差分方法则更容易产生局部振荡误差。

因此，在二维波动方程与量子波函数模拟问题中，FFT 谱方法通常具有更高数值效率与更优传播质量。


\subsection{结论}


本文从二维含势场薛定谔方程出发，推导了其傅里叶谱方法离散格式。核心思想是利用傅里叶变换的微分性质，将拉普拉斯算符转化为频域中的代数乘法，从而高效计算哈密顿算符。再结合显式 Euler 时间推进，可以得到波函数的逐步演化公式：
\[
\psi^{n+1}
=
\psi^n
-i\Delta t
\left[
-\Delta\psi^n+V\psi^n
\right]
\]

其中
\[
\Delta\psi^n
=
\mathcal{F}^{-1}
\left[
-(k_x^2+k_y^2)\mathcal{F}[\psi^n]
\right].
\]

综合实验结果可以发现：
有限差分方法实现简单；FFT 谱方法精度更高；
FFT 方法能够更准确描述波函数传播与干涉；
在大规模模拟中，FFT 方法具有更高效率。

\begin{thebibliography}{99}

\bibitem{griffiths}
D. J. Griffiths,
\textit{Introduction to Quantum Mechanics},
2nd Edition,
Pearson Prentice Hall, 2005.

\bibitem{sakurai}
J. J. Sakurai and J. Napolitano,
\textit{Modern Quantum Mechanics},
2nd Edition,
Cambridge University Press, 2017.

\bibitem{trefethen}
L. N. Trefethen,
\textit{Spectral Methods in MATLAB},
SIAM, 2000.

\bibitem{boyd}
J. P. Boyd,
\textit{Chebyshev and Fourier Spectral Methods},
2nd Edition,
Dover Publications, 2001.

\bibitem{strang}
G. Strang,
\textit{Computational Science and Engineering},
Wellesley-Cambridge Press, 2007.

\bibitem{press}
W. H. Press, S. A. Teukolsky, 
W. T. Vetterling, and B. P. Flannery,
\textit{Numerical Recipes: The Art of Scientific Computing},
3rd Edition,
Cambridge University Press, 2007.

\bibitem{fft}
J. W. Cooley and J. W. Tukey,
``An Algorithm for the Machine Calculation of Complex Fourier Series,''
\textit{Mathematics of Computation},
vol. 19, no. 90, pp. 297--301, 1965.

\end{thebibliography}

\end{document}
